% ============================================================================
% Fichier  : partie4/P04_C13_cartographie_traitement.tex
% Titre    : Cartographie et Traitement des Risques
% Auteur   : MINKA MI NGUIDJOÏ Thierry Emmanuel
% ============================================================================
% Sources : ISO/IEC 27005:2022 ; ISO 31000:2018. Mairie d'Eséka : cas fictif.
% ============================================================================

\chapter{Cartographie et Traitement des Risques}
\label{chap:traitement}
\chaptermark{Cartographie et traitement}

\epigraph{%
    « Cartographier, c'est transformer une inquiétude diffuse\\
    en une liste de décisions à prendre. »%
}{Principe du cours}

\niveauChapitre{Intermédiaire}{Les Chapitres~\ref{chap:vocabrisque} et~\ref{chap:methodes}.}

\begin{boxobjectifs}
\begin{itemize}
    \item Construire et lire une matrice des risques.
    \item Hiérarchiser les risques par niveau.
    \item Énoncer les quatre stratégies de traitement.
    \item Choisir une stratégie selon le risque et le coût.
    \item Formaliser le risque résiduel et son acceptation.
\end{itemize}
\end{boxobjectifs}

\bigskip

% ============================================================================
\section*{Question directrice : une fois les risques connus, que décide-t-on ?}
% ============================================================================

Connaître ses risques ne sert à rien si l'on n'en fait rien. Ce chapitre
transforme la liste des risques en décisions~: les classer, puis, pour chacun,
choisir un traitement.

\medskip
\textbf{Plan du chapitre.}
\begin{enumerate}
    \item La matrice des risques (section~\ref{sec:c13-matrice}).
    \item Hiérarchiser (section~\ref{sec:c13-hierarchie}).
    \item Les quatre stratégies (section~\ref{sec:c13-strategies}).
    \item Choisir un traitement (section~\ref{sec:c13-choisir}).
    \item Le risque résiduel (section~\ref{sec:c13-residuel}).
\end{enumerate}

\separateur

% ============================================================================
\section{La Matrice des Risques}
\label{sec:c13-matrice}
% ============================================================================

\begin{boxdef}
La \termecle{matrice des risques} croise la \emph{probabilité} et l'\emph{impact}
pour situer chaque risque dans une grille et en déduire un niveau.
\end{boxdef}

\begin{tableaucours}{L{2.4cm}L{2.0cm}L{2.0cm}L{2.0cm}L{2.0cm}}{%
    \tableauHeader{Impact / Prob.} &
    \tableauHeader{Faible} &
    \tableauHeader{Moyenne} &
    \tableauHeader{Élevée} &
    \tableauHeader{Très él.}
}
Mineur & Faible & Faible & Moyen & Moyen \\
Modéré & Faible & Moyen & Moyen & Élevé \\
Majeur & Moyen & Moyen & Élevé & Très él. \\
Critique & Moyen & Élevé & Très él. & Très él. \\
\end{tableaucours}
\vspace{-0.8em}
\begin{center}{\small\itshape Tableau~13.1. Matrice probabilité $\times$ impact}\end{center}

% ============================================================================
\section{Hiérarchiser}
\label{sec:c13-hierarchie}
% ============================================================================

La matrice hiérarchise~: on traite d'abord les risques « très élevés » et
« élevés », on surveille les « moyens », on documente l'acceptation des
« faibles ». Cette hiérarchie fonde le \termecle{plan de traitement}~: on ne
peut pas tout traiter en même temps, donc on priorise.

\begin{boxprocedure}
\textbf{Constatation.} À Eséka, le risque « perte de l'état civil » est très
élevé~; le risque « défiguration du site web » est moyen.

\medskip
\textbf{Analyse.} La priorité va à l'état civil~: impact critique, probabilité
démontrée. Le site web, moins critique, sera traité ensuite. Hiérarchiser,
c'est choisir où mettre ses moyens limités.
\end{boxprocedure}

% ============================================================================
\section{Les Quatre Stratégies}
\label{sec:c13-strategies}
% ============================================================================

\begin{boxdef}
Pour chaque risque, quatre \termecle{stratégies de traitement}~:
\begin{itemize}
    \item \termecle{Éviter}~: supprimer la cause (renoncer à l'activité
          risquée).
    \item \termecle{Réduire}~: diminuer la probabilité ou l'impact par des
          mesures.
    \item \termecle{Transférer}~: reporter le risque (assurance,
          sous-traitance).
    \item \termecle{Accepter}~: assumer le risque résiduel en connaissance de
          cause.
\end{itemize}
\end{boxdef}

% ============================================================================
\section{Choisir un Traitement}
\label{sec:c13-choisir}
% ============================================================================

\begin{tableaucours}{L{4.2cm}L{3.0cm}L{4.4cm}}{%
    \tableauHeader{Risque} &
    \tableauHeader{Stratégie} &
    \tableauHeader{Mesure}
}
Perte de l'état civil & Réduire & Sauvegarde 3-2-1, test de restauration \\
Fuite du fichier taxes & Réduire & Chiffrement, accès nominatif \\
Panne électrique & Transférer & Contrat de maintenance, onduleur loué \\
Site web décoratif & Accepter & Risque résiduel faible, documenté \\
\end{tableaucours}
\vspace{-0.8em}
\begin{center}{\small\itshape Tableau~13.2. Choix de traitement à la mairie}\end{center}

Le choix met en balance le \emph{coût} du traitement et le \emph{gain} de
réduction du risque~: réduire un risque faible à grands frais est un mauvais
arbitrage.

\begin{boxmetier}{Direction générale}
\textbf{Accepter est une décision, pas un oubli.} Accepter un risque est
légitime, à condition que ce soit explicite, documenté et décidé par la
direction. Un risque « accepté » par simple inaction n'est pas accepté~: il est
ignoré.
\end{boxmetier}

% ============================================================================
\section{Le Risque Résiduel}
\label{sec:c13-residuel}
% ============================================================================

Après traitement, il reste un \termecle{risque résiduel}. La norme exige que la
direction l'\termecle{accepte formellement} (Chapitre~\ref{chap:dda}). Cette
acceptation est tracée~: c'est la direction qui porte, en connaissance de
cause, ce qui n'a pas été éliminé.

\separateur

\begin{boxresume}
\textbf{Ce qu'il faut retenir de ce chapitre~:}
\begin{itemize}
    \item La \textbf{matrice} situe chaque risque et en déduit un niveau.
    \item On \textbf{hiérarchise} pour traiter d'abord les risques les plus
          forts.
    \item Quatre \textbf{stratégies}~: éviter, réduire, transférer, accepter.
    \item On \textbf{choisit} en pesant coût et gain.
    \item Le \textbf{risque résiduel} est accepté formellement par la
          direction.
\end{itemize}
\end{boxresume}

\begin{boxinfo}
\textbf{Points de contrôle de l'élève.} Je sais~: lire une matrice des
risques~; hiérarchiser~; énoncer les quatre stratégies~; choisir un
traitement~; expliquer l'acceptation du risque résiduel.
\end{boxinfo}

\bigskip

\begin{boxexo}[title={\ding{45}\quad Exercice~13.1. Lire la matrice \niveaubase}]
Un risque de probabilité élevée et d'impact critique~: quel niveau~?
\end{boxexo}

\begin{boxexo}[title={\ding{45}\quad Exercice~13.2. Les stratégies \niveaubase}]
Citez les quatre stratégies de traitement du risque.
\end{boxexo}

\begin{boxexo}[title={\ding{45}\quad Exercice~13.3. Hiérarchiser \niveaubase}]
Entre un risque « très élevé » et un « moyen », lequel traite-t-on d'abord~?
\end{boxexo}

\begin{boxexo}[title={\ding{45}\quad Exercice~13.4. Associer \niveaumoyen}]
Associez une stratégie à chacun~: perte de l'état civil~; panne électrique~;
site web décoratif.
\end{boxexo}

\begin{boxexo}[title={\ding{45}\quad Exercice~13.5. Coût-gain \niveaumoyen}]
Pourquoi réduire un risque faible à grands frais est-il un mauvais arbitrage~?
\end{boxexo}

\begin{boxexo}[title={\ding{45}\quad Exercice~13.6. Transférer \niveaumoyen}]
Donnez un exemple de risque que la mairie peut transférer, et comment.
\end{boxexo}

\begin{boxexo}[title={\ding{45}\quad Exercice~13.7. Accepter ou ignorer \niveaumoyen}]
Quelle différence entre un risque accepté et un risque ignoré~?
\end{boxexo}

\begin{boxexo}[title={\ding{45}\quad Exercice~13.8. Résiduel \niveaumoyen}]
Qui accepte le risque résiduel, et pourquoi doit-ce être tracé~?
\end{boxexo}

\begin{boxexo}[title={\ding{45}\quad Exercice~13.9. Cartographie \niveauexpert}]
Dressez la cartographie de cinq risques de la mairie dans la matrice, avec leur
niveau.
\end{boxexo}

\begin{boxexo}[title={\ding{45}\quad Exercice~13.10. Plan de traitement \niveauexpert}]
Construisez un plan de traitement pour ces cinq risques (stratégie, mesure,
responsable, échéance).
\end{boxexo}

\bigskip

\begin{boxinfo}
\textbf{Pour aller plus loin.}
\begin{itemize}
    \item ISO/IEC 27005~\cite{iso27005}~; ISO 31000~\cite{iso31000}.
\end{itemize}
\end{boxinfo}

\begin{boxinfo}
\textbf{Transition.} Le plan de traitement se formalise dans un document
central~: le Chapitre~\ref{chap:dda} présente la déclaration d'applicabilité.
\end{boxinfo}
